
% https://akaedu.github.io/book/apa.html

\chapter{字符编码}
\BookAnchor{db.apa}\BookAnchor{db.apa.id2906420}%


\BookNav{\hyperlink{db.ch37s05}{上一页}\enspace\textperiodcentered\enspace \hyperlink{db.index}{上一级}\enspace\textperiodcentered\enspace \hyperlink{db.index}{起始页}\enspace\textperiodcentered\enspace \hyperlink{db.apas01}{下一页}}


% https://akaedu.github.io/book/apas01.html

\section{ASCII码}
\BookAnchor{db.apas01}\BookAnchor{db.apas01.id2906424}%
ASCII码的取值范围是0\textasciitilde{}\BookSplit{}127，可以用7个bit表示。C语言中\BookCode{char}型变量的大小规定为一字节，如果存放ASCII码则只用到低7位，高位为0。以下是ASCII码表：


\BookFigure{ASCII码表}{\BookImage{images/app-encoding.ascii.png}{0.891}}{\BookAnchor{db.apas01.id2905808}}


绝大多数计算机的一个字节是8位，取值范围是0\textasciitilde{}\BookSplit{}255，而ASCII码并没有规定编号为128\textasciitilde{}\BookSplit{}255的字符，为了能表示更多字符，各厂商制定了很多种ASCII码的扩展规范。注意，虽然通常把这些规范称为扩展ASCII码（Extended ASCII）\BookAnchor{db.apas01.id2905832}，但其实它们并不属于ASCII码标准。例如以下这种扩展ASCII码由IBM制定，在字符终端下被广泛采用，其中包含了很多表格边线字符用来画界面。


\BookFigure{IBM的扩展ASCII码表}{\BookImage{images/app-encoding.extascii.png}{0.861}}{\BookAnchor{db.apas01.id2906450}}


在图形界面中最广泛使用的扩展ASCII码是ISO-8859-1，也称为Latin-1，其中包含欧洲各国语言中最常用的非英文字母，但毕竟只有128个字符，某些语言中的某些字母没有包含。如下表所示。


\BookFigure{ISO-8859-1}{\BookImage{images/app-encoding.latin1.png}{0.624}}{\BookAnchor{db.apas01.id2906470}}


编号为128\textasciitilde{}\BookSplit{}159的是一些控制字符，在上表中没有列出。

\BookNav{\hyperlink{db.apa}{上一页}\enspace\textperiodcentered\enspace \hyperlink{db.apa}{上一级}\enspace\textperiodcentered\enspace \hyperlink{db.index}{起始页}\enspace\textperiodcentered\enspace \hyperlink{db.apas02}{下一页}}


% https://akaedu.github.io/book/apas02.html

\section{Unicode和UTF-8}
\BookAnchor{db.apas02}\BookAnchor{db.apas02.id2906489}%
为了统一全世界各国语言文字和专业领域符号（例如数学符号、乐谱符号）的编码，ISO制定了ISO 10646\BookAnchor{db.apas02.id2906502}标准，也称为UCS（Universal Character Set）\BookAnchor{db.apas02.id2906510}。UCS编码的长度是31位，可以表示2\textsuperscript{31}个字符。如果两个字符编码的高位相同，只有低16位不同，则它们属于一个平面（Plane）\BookAnchor{db.apas02.id2906524}，所以一个平面由2\textsuperscript{16}个字符组成。目前常用的大部分字符都位于第一个平面（编码范围是U-00000000\textasciitilde{}\BookSplit{}U-0000FFFD），称为BMP（Basic Multilingual Plane）\BookAnchor{db.apas02.id2906537}或Plane 0，为了向后兼容，其中编号为0\textasciitilde{}\BookSplit{}256的字符和Latin-1相同。UCS编码通常用U-xxxxxxxx这种形式表示，而BMP的编码通常用U+\BookSplit{}xxxx这种形式表示，其中x是十六进制数字。在ISO制定UCS的同时，另一个由厂商联合组织也在着手制定这样的编码，称为Unicode\BookAnchor{db.apas02.id2906555}，后来两家联手制定统一的编码，但各自发布各自的标准文档，所以UCS编码和Unicode码是相同的。



有了字符编码，另一个问题就是这样的编码在计算机中怎么表示。现在已经不可能用一个字节表示一个字符了，最直接的想法就是用四个字节表示一个字符，这种表示方法称为UCS-4\BookAnchor{db.apas02.id2906571}或UTF-32\BookAnchor{db.apas02.id2906578}，UTF是Unicode Transformation Format\BookAnchor{db.apas02.id2906586}的缩写。一方面这样比较浪费存储空间，由于常用字符都集中在BMP，高位的两个字节通常是0，如果只用ASCII码或Latin-1，高位的三个字节都是0。另一种比较节省存储空间的办法是用两个字节表示一个字符，称为UCS-2\BookAnchor{db.apas02.id2906598}或UTF-16\BookAnchor{db.apas02.id2906605}，这样只能表示BMP中的字符，但BMP中有一些扩展字符，可以用两个这样的扩展字符表示其它平面的字符，称为Surrogate Pair\BookAnchor{db.apas02.id2906615}。无论是UTF-32还是UTF-16都有一个更严重的问题是和C语言不兼容，在C语言中0字节表示字符串结尾，库函数\BookCode{strlen}、\BookCode{strcpy}等等都依赖于这一点，如果字符串用UTF-32存储，其中有很多0字节并不表示字符串结尾，这就乱套了。



UNIX之父Ken Thompson提出的UTF-8\BookAnchor{db.apas02.id2906642}编码很好地解决了这些问题，现在得到广泛应用。UTF-8具有以下性质：


\begin{itemize}

\item 编码为U+\BookSplit{}0000\textasciitilde{}\BookSplit{}U+\BookSplit{}007F的字符只占一个字节，就是0x00\textasciitilde{}\BookSplit{}0x7F，和ASCII码兼容。

\item 编码大于U+\BookSplit{}007F的字符用2\textasciitilde{}\BookSplit{}6个字节表示，每个字节的最高位都是1，而ASCII码的最高位都是0，因此非ASCII码字符的表示中不会出现ASCII码字节（也就不会出现0字节）。

\item 用于表示非ASCII码字符的多字节序列中，第一个字节的取值范围是0xC0\textasciitilde{}\BookSplit{}0xFD，根据它可以判断后面有多少个字节也属于当前字符的编码。后面每个字节的取值范围都是0x80\textasciitilde{}\BookSplit{}0xBF，见下面的详细说明。

\item UCS定义的所有2\textsuperscript{31}个字符都可以用UTF-8编码表示出来。

\item UTF-8编码最长6个字节，BMP字符的UTF-8编码最长三个字节。

\item 0xFE和0xFF这两个字节在UTF-8编码中不会出现。
\end{itemize}


具体来说，UTF-8编码有以下几种格式：


\begin{BookLiteral}
\BookLiteralLine{U-00000000 – U-0000007F: 0xxxxxxx}
\BookLiteralLine{U-00000080 – U-000007FF: 110xxxxx 10xxxxxx}
\BookLiteralLine{U-00000800 – U-0000FFFF: 1110xxxx 10xxxxxx 10xxxxxx}
\BookLiteralLine{U-00010000 – U-001FFFFF: 11110xxx 10xxxxxx 10xxxxxx 10xxxxxx}
\BookLiteralLine{U-00200000 – U-03FFFFFF: 111110xx 10xxxxxx 10xxxxxx 10xxxxxx 10xxxxxx}
\BookLiteralLine{U-04000000 – U-7FFFFFFF: 1111110x 10xxxxxx 10xxxxxx 10xxxxxx 10xxxxxx 10xxxxxx}
\end{BookLiteral}


第一个字节要么最高位是0（ASCII字节），要么最高两位都是1，最高位之后1的个数决定后面有多少个字节也属于当前字符编码，例如111110xx，最高位之后还有四个1，表示后面有四个字节也属于当前字符的编码。后面每个字节的最高两位都是10，可以和第一个字节区分开。这样的设计有利于误码同步，例如在网络传输过程中丢失了几个字节，很容易判断当前字符是不完整的，也很容易找到下一个字符从哪里开始，结果顶多丢掉一两个字符，而不会导致后面的编码解释全部混乱了。上面的格式中标为x的位就是UCS编码，最后一种6字节的格式中x位有31个，可以表示31位的UCS编码，UTF-8就像一列火车，第一个字节是车头，后面每个字节是车厢，其中承载的货物是UCS编码。UTF-8规定承载的UCS编码以大端表示，也就是说第一个字节中的x是UCS编码的高位，后面字节中的x是UCS编码的低位。



例如U+\BookSplit{}00A9（©字符）的二进制是10101001，编码成UTF-8是11000010 10101001（0xC2 0xA9），但不能编码成11100000 10000010 10101001，UTF-8规定每个字符只能用尽可能少的字节来编码。

\BookNav{\hyperlink{db.apas01}{上一页}\enspace\textperiodcentered\enspace \hyperlink{db.apa}{上一级}\enspace\textperiodcentered\enspace \hyperlink{db.index}{起始页}\enspace\textperiodcentered\enspace \hyperlink{db.apas03}{下一页}}


% https://akaedu.github.io/book/apas03.html

\section{在Linux C编程中使用Unicode和UTF-8}
\BookAnchor{db.apas03}\BookAnchor{db.apas03.id2906750}%
目前各种Linux发行版都支持UTF-8编码，当前系统的语言和字符编码设置保存在一些环境变量中，可以通过\BookCode{locale}命令查看：


\BookTerminal{code/apas03/01.txt}{9}{15}


常用汉字也都位于BMP中，所以一个汉字的存储通常占3个字节。例如编辑一个C程序：


\BookSource{numbers=none}{code/apas03/02.pv}{9}{7}


源文件是以UTF-8编码存储的：


\BookTerminal{code/apas03/03.txt}{9}{7}


其中八进制的\BookCode{344 375 240}（十六进制\BookCode{e4 bd a0}）就是“你”的UTF-8编码，八进制的\BookCode{345 245 275}（十六进制\BookCode{e5 a5 bd}）就是“好”。把它编译成目标文件，\BookCode{\BookQuoteChar{}你好\textbackslash{}n\BookQuoteChar{}}这个字符串就成了这样一串字节：\BookCode{e4 bd a0 e5 a5\BookSplit{} bd 0a 00}，汉字在其中仍然是UTF-8编码的，一个汉字占3个字节，这种字符在C语言中称为多字节字符（Multibyte Character）\BookAnchor{db.apas03.id2906864}。运行这个程序相当于把这一串字节\BookCode{write}到当前终端的设备文件。如果当前终端的驱动程序能够识别UTF-8编码就能打印出汉字，如果当前终端的驱动程序不能识别UTF-8编码（比如一般的字符终端）就打印不出汉字。也就是说，像这种程序，识别汉字的工作既不是由C编译器做的也不是由\BookCode{libc}做的，C编译器原封不动地把源文件中的UTF-8编码复制到目标文件中，\BookCode{libc}只是当作以0结尾的字符串原封不动地\BookCode{write}给内核，识别汉字的工作是由终端的驱动程序做的。



但是仅有这种程度的汉字支持是不够的，有时候我们需要在C程序中操作字符串里的字符，比如求字符串\BookCode{\BookQuoteChar{}你好\textbackslash{}n\BookQuoteChar{}}中有几个汉字或字符，用\BookCode{strlen}就不灵了，因为\BookCode{strlen}只看结尾的0字节而不管字符串里存的是什么，求出来的是字节数7。为了在程序中操作Unicode字符，C语言定义了宽字符（Wide Character）\BookAnchor{db.apas03.id2906938}类型\BookCode{wchar\_\BookSplit{}t}和一些库函数。在字符常量或字符串字面值前面加一个L就表示宽字符常量或宽字符串，例如定义\BookCode{wchar\_\BookSplit{}t c =\BookSplit{} L'你';\BookSplit{}}，变量\BookCode{c}的值就是汉字“你”的31位UCS编码，而\BookCode{L\BookQuoteChar{}你好\textbackslash{}n\BookQuoteChar{}}就相当于\BookCode{\{L'你',\BookSplit{} L'好',\BookSplit{} L'\textbackslash{}n',\BookSplit{} 0\}}，\BookCode{wcslen}函数就可以取宽字符串中的字符个数。看下面的程序：


\BookSource{numbers=left,xleftmargin=19pt}{code/apas03/04.pv}{9}{13}


宽字符串\BookCode{L\BookQuoteChar{}你好\textbackslash{}n\BookQuoteChar{}}在源代码中当然还是存成UTF-8编码的，但编译器会把它变成4个UCS编码\BookCode{0x00004f60 0x0\BookSplit{}000597d 0x0000\BookSplit{}000a 0x0000000\BookSplit{}0}保存在目标文件中，按小端存储就是\BookCode{60 4f 00 00 7d\BookSplit{} 59 00 00 0a 0\BookSplit{}0 00 00 00 00 \BookSplit{}00 00}，用\BookCode{od}命令查看目标文件应该能找到这些字节。


\BookTerminal{code/apas03/05.txt}{4}{2}


\BookCode{printf}的\BookCode{\%ls}转换说明表示把后面的参数按宽字符串解释，不是见到0字节就结束，而是见到UCS编码为0的字符才结束，但是要\BookCode{write}到终端仍然需要以多字节编码输出，这样终端驱动程序才能识别，所以\BookCode{printf}在内部把宽字符串转换成多字节字符串再\BookCode{write}出去。事实上，C标准并没有规定多字节字符必须以UTF-8编码，也可以使用其它的多字节编码，在运行时根据环境变量确定当前系统的编码，所以在程序开头需要调用\BookCode{setlocale}获取当前系统的编码设置，如果当前系统是UTF-8的，\BookCode{printf}就把UCS编码转换成UTF-8编码的多字节字符串再\BookCode{write}出去。一般来说，程序在做内部计算时通常以宽字符编码，如果要存盘或者输出给别的程序，或者通过网络发给别的程序，则采用多字节编码。



关于Unicode和UTF-8本节只介绍了最基本的概念，部分内容出自\hyperlink{db.bi01.bibli.unicodefaq}{[Unicode FAQ]}，读者可进一步参考这篇文章。

\BookNav{\hyperlink{db.apas02}{上一页}\enspace\textperiodcentered\enspace \hyperlink{db.apa}{上一级}\enspace\textperiodcentered\enspace \hyperlink{db.index}{起始页}\enspace\textperiodcentered\enspace \hyperlink{db.apb}{下一页}}
